% To be reformatted in the journal style upon acceptance.
\documentclass[11pt]{amsart}

\usepackage{amsmath,amssymb}
\usepackage[hidelinks]{hyperref}

\newtheorem{thm}{Theorem}[section]
\newtheorem{lem}[thm]{Lemma}
\newtheorem{cor}[thm]{Corollary}
\newtheorem{prop}[thm]{Proposition}
\newtheorem{conj}[thm]{Conjecture}
\theoremstyle{definition}
\newtheorem{defi}[thm]{Definition}
\newtheorem{rem}[thm]{Remark}
\newtheorem{ques}[thm]{Question}

\newcommand{\I}{\mathbb{I}}
\newcommand{\cube}{\square_{\mathrm{DM}}}
\newcommand{\cubeK}{\square_{\mathrm{KL}}}
\newcommand{\el}{\mathrm{el}}
\newcommand{\DM}{\mathrm{DM}}
\newcommand{\Wtest}{\mathsf{W}_{\mathrm{test}}}
\newcommand{\Wtype}{\mathsf{W}_{\mathrm{type}}}
\newcommand{\Wmin}{\mathsf{W}_{\mathrm{min}}}
\newcommand{\Stab}{\mathrm{Stab}}
\newcommand{\sw}{\mathsf{sw}}
\newcommand{\nb}{\mathsf{n}}
\newcommand{\gtw}{\mathsf{g}}
\newcommand{\Kl}{K}
\newcommand{\hopush}{\mathrm{hocolim}}

\begin{document}

\title[Isotropy quotients and the test comparison]{The join of two
  projective spaces:\\ isotropy quotients and the test comparison\\
  for De Morgan cubical sets}

\author{Ryota Shibusawa}
\address{Daiichi Institute of Technology, Japan}
\email{r-shibusawa@daiichi-koudai.ac.jp}

\subjclass[2020]{55U35; 18N40, 55U10, 06D30, 03B38}
\keywords{cubical sets, cubical type theory, model structure, test
  category, localizer, orbit space, De Morgan algebra, join}

\begin{abstract}
  A companion paper re-opened the question whether the cubical-type
  model structure on De Morgan cubical sets presents spaces.  We
  first prove the upper bound: on any totally aspherical test cube
  category the cubical-type weak equivalences are contained in the
  test weak equivalences, by a one-paragraph localizer argument; on
  cartesian cubes the inclusion is strict, by Coquand's symmetric
  square.  We then test sharpness on the isotropy quotients of the
  De Morgan square.  Quotients by a cyclic group whose generator
  fixes a cell contract in both structures, by an explicit
  equivariant interpolation that connections make available;
  quotients by a freely acting group are $K(H,1)$ in both
  structures.  For the two remaining groups --- the Klein group
  generated by the transposition and the double negation, and the
  full hyperoctahedral group --- the category of elements of the
  quotient is computed exactly: for the Klein group it is the
  \emph{join} $\mathbb{RP}^\infty \star \mathbb{RP}^\infty$, simply
  connected with third homology $\mathbb{Z}/2$, so the quotient is
  contractible in neither structure; a machine reachability census
  over twisted-symmetric cells gives an independent syntactic proof.
  Consequently every known candidate witness agrees, and no
  automorphism quotient of the square separates the two localizers.
  We conjecture that the agreement is general, and isolate the
  presheaf-level obstacle: the geometric realization of these
  quotients is contractible, and every stabilized cell is a face of
  a free cell, so neither cellular computation nor na\"ive
  stratification sees the answer.
\end{abstract}

\maketitle

\section{Introduction}
\label{sec:intro}

Let $\cube$ be the De Morgan cube category --- the Lawvere theory of
De Morgan algebras, restricted to the free algebras
$\DM(n)$ on the interval generators --- and let
$\widehat{\cube}$ be its presheaf category of \emph{De Morgan
  cubical sets}, the site of the CCHM model of cubical type
theory \cite{CCHM}.  Two Quillen model structures live on
$\widehat{\cube}$ with the same cofibrations (all monomorphisms):
the \emph{cubical-type} model structure, whose fibrations are the
uniform Kan fibrations of cubical type theory, with weak
equivalences $\Wtype$; and the \emph{test} model structure of
Cisinski's theory of test categories \cite{Cisinski}, whose weak
equivalences $\Wtest$ are the maps $f$ such that $\el(f)$ is a weak
equivalence of categories, where $\el$ denotes the category of
elements.  Since $\cube$ is a strict test category
\cite{BuchholtzMorehouse, Maltsiniotis}, the test structure
presents the homotopy theory of spaces.  Whether the type structure
does --- that is, whether $\Wtype = \Wtest$ --- is the
\emph{test comparison problem} for cubical type theory.

For cartesian cubes the answer is negative: Coquand's symmetric
square $Q$ is test-trivial but not type-contractible
\cite{CoqNote, ACCRS}.  For De Morgan cubes the folklore witness
was the reversal quotient $L = \I/(x \sim \neg x)$; a companion
paper \cite{PaperK} proved $L$ is a $K(\mathbb{Z}/2,1)$ in the type
structure, and its sequel \cite{PaperL} showed the same holds on
the test side ($\el(L) \simeq B\mathbb{Z}/2$, by freeness of the
reversal action on cells), so the witness is \emph{void} and the
De Morgan comparison problem is open.

This paper reports three groups of results.

\subsection*{The upper bound}
Theorem~\ref{thm:upper}: for every totally aspherical test cube
category whose type structure is generated by
pushout products of endpoint inclusions with monomorphisms,
\[
  \Wtype \subseteq \Wtest .
\]
The proof is a one-paragraph localizer argument.  Combined with
Coquand's square, it shows the inclusion is \emph{strict} for
cartesian cubes; and its contrapositive --- test-nontrivial implies
type-nontrivial --- becomes the engine of the second half of the
paper.

\subsection*{Isotropy quotients}
The remaining question is the lower bound, equivalently: is every
test-trivial map type-trivial?  A separating witness must be
test-trivial and type-nontrivial.  All known candidates are
quotients $\square^n/H$ by subgroups $H$ of the hyperoctahedral
automorphism group \cite{PaperL}, so we compute both sides for
\emph{all} subgroups of $\mathrm{Aut}(\square^2) = G_2$
(dihedral of order $8$), and for a representative
three-dimensional family:
\begin{itemize}
\item \emph{Fixed-stratum quotients} (Theorem~\ref{thm:retract}):
  if a cyclic group's generator fixes at least one cell, the
  quotient is contractible in the type structure --- an explicit
  equivariant interpolation, available because connections let one
  solve the equivariance constraints around the cycle --- hence
  also in the test structure, by Theorem~\ref{thm:upper}.  This
  uniformly contracts the symmetric square $\square^2/\langle
  \sw\rangle$, the twisted square $\square^2/\langle
  \gtw\rangle$, and a $\mathbb{Z}/3$ quotient of the cube.
\item \emph{Free quotients} (Theorem~\ref{thm:free}): if $H$ acts
  freely on cells, the quotient is a $K(H,1)$ in the type
  structure, matching $\el(\square^n/H) \simeq BH$ \cite{PaperL}
  on the nose.
\item \emph{Mixed quotients} (Theorem~\ref{thm:join}): for the
  Klein group $\Kl = \{1, \sw, \nb, \gtw\}$ (transposition, double
  negation, and their product) no cell is fixed by all of $\Kl$,
  and the strict quotient $\el(\square^2/\Kl)$ is computed exactly:
  it is the join
  \[
    \el(\square^2/\Kl) \;\simeq\; B\mathbb{Z}/2 \star B\mathbb{Z}/2
    \;=\; \mathbb{RP}^\infty \star \mathbb{RP}^\infty ,
  \]
  simply connected with $H_3 = \mathbb{Z}/2$: contractible in
  \emph{neither} structure.  The same method computes the full
  hyperoctahedral quotient.  A twisted reachability census ---
  $72$ symmetric squares, of which $66$ do contract and the
  generic one is isolated --- gives an independent, purely
  syntactic proof of type-noncontractibility in the style of
  \cite{CoqNote, PaperK}.
\end{itemize}

\subsection*{Agreement}
Every candidate tested --- the reversal circle, the symmetric and
twisted squares, all free quotients, the Klein and hyperoctahedral
quotients, the $\mathbb{Z}/3$ cube --- has the \emph{same} answer
in both structures (Table~\ref{tab:agree}).  In particular no
automorphism quotient of the square separates the localizers
(Corollary~\ref{cor:nosep}).  We conjecture the agreement is
general (Conjecture~\ref{conj:agree}), identify why the mixed-case
type computation is genuinely hard (the realization is
contractible and every stabilized cell is a face of a free cell,
Remark~\ref{rem:obstacle}), and record along the way that the De
Morgan cube category is \emph{not} an Eilenberg--Zilber category
(Proposition~\ref{prop:ez}) --- the Kleene midpoint has no
(split epi, mono) factorization --- which is why elegance-style
techniques are unavailable on this site.

All machine verifications are reproduced by
\texttt{scripts/isotropy.py} in the artifact \cite{Artifact}.

\section{Two model structures}
\label{sec:prelim}

We write $\DM(n)$ for the free De Morgan algebra on $x_1, \dots,
x_n$; $\cube$ has objects $[n]$ and $\cube([m],[n]) =
\DM(m)^n$.  Recall $|\DM(1)| = 6$, $|\DM(2)| = 168$,
$|\DM(3)| = 7{,}828{,}354$; elements are monotone functions on the
$2n$-coordinate \emph{literal cube}, with $\neg$ acting by
complement-and-swap within literal pairs \cite{PaperG}.

The \emph{cubical-type model structure} on $\widehat\cube$ has
cofibrations the monomorphisms and fibrations the maps with the
right lifting property against all pushout products
$\delta_e \mathbin{\hat\otimes} m$ of an endpoint inclusion
$\delta_e : 1 \to \I$ ($e \in \{0,1\}$) with a monomorphism $m$;
this is the Cisinski--Sattler construction underlying the CCHM
model \cite{CCHM, Sattler}.  Its weak equivalences are $\Wtype$.
The \emph{test model structure} has the same cofibrations and weak
equivalences $\Wtest = \{f : N\el(f) \text{ a simplicial weak
  equivalence}\}$; it exists and presents spaces because $\cube$
is a strict --- indeed totally aspherical --- test category
\cite{Cisinski, Maltsiniotis, BuchholtzMorehouse}.

An \emph{$A$-localizer} \cite[\S1.4]{Cisinski} is a class $\mathsf
W$ of maps of $\widehat A$ containing the trivial fibrations,
satisfying two-out-of-three, and whose monomorphic members are
closed under pushout, transfinite composition, and retract.
$\Wtest$ is an $A$-localizer, and $\Wtype$ is one as well
\cite{PaperL}; $\Wmin$ denotes the minimal localizer.

$\mathrm{Aut}_{\cube}([n]) = G_n$, the hyperoctahedral group
$\Sigma_n \ltimes (\mathbb{Z}/2)^n$, acting on cells
$(a_1, \dots, a_n) \in \DM(m)^n$ by permuting coordinates and
applying $\neg$ \cite{PaperL}.  For $n = 2$ we write
$\sw(a,b) = (b,a)$, $\nb(a,b) = (\neg a, \neg b)$,
$\gtw = \sw\nb : (a,b) \mapsto (\neg b, \neg a)$, and
$\nb_1(a,b) = (\neg a, b)$; $G_2 = \langle \sw, \nb_1\rangle$ is
dihedral of order $8$ and $\Kl = \{1, \sw, \nb, \gtw\}$ is a Klein
four-group.

\section{The upper bound}
\label{sec:upper}

\begin{thm}
  \label{thm:upper}
  Let $A$ be a totally aspherical test cube category (De Morgan,
  Kleene, Boolean, or cartesian) whose cubical-type model structure
  is cofibrantly generated with generating trivial cofibrations
  $\{\delta_e \mathbin{\hat\otimes} m : m \text{ mono},\,
  e \in \{0,1\}\}$.  Then $\Wtype \subseteq \Wtest$.
\end{thm}

\begin{proof}
  $\Wtest$ is an $A$-localizer.  (i)~Since $A$ is totally
  aspherical, $\el$ preserves finite products up to weak
  equivalence \cite[\S1.6]{Cisinski}, and $\el(\I)$ has a terminal
  object; hence the projection $X \times \I \to X$ is in $\Wtest$
  for every $X$, and by two-out-of-three so is each end inclusion
  $X \times \{e\} \hookrightarrow X \times \I$ --- a monomorphism.
  (ii)~For a monomorphism $m : X \rightarrowtail Y$, the map
  $\delta_e \mathbin{\hat\otimes} m$ is the comparison
  $Y \cup_X (X \times \I) \to Y \times \I$.  The inclusion
  $Y \hookrightarrow Y \cup_X (X \times \I)$ is a pushout of the
  monomorphic test equivalence $X \hookrightarrow X \times \I$,
  hence in $\Wtest$ by the localizer axioms; the end inclusion
  $Y \hookrightarrow Y \times \I$ is in $\Wtest$ by (i); so
  $\delta_e \mathbin{\hat\otimes} m \in \Wtest$ by
  two-out-of-three.  (iii)~The trivial cofibrations of the type
  structure are retracts of transfinite compositions of pushouts of
  the generators, all monomorphisms, so they lie in $\Wtest$; the
  trivial fibrations lie in every localizer; and an arbitrary map
  of $\Wtype$ factors as a trivial cofibration followed by a
  trivial fibration.
\end{proof}

\begin{cor}
  \label{cor:cartesian}
  On cartesian cubical sets $\Wtype \subsetneq \Wtest$: the
  inclusion holds by Theorem~\ref{thm:upper}, and it is strict
  because Coquand's symmetric square is test-trivial
  \cite{ACCRS} but not type-trivial \cite{CoqNote}.
\end{cor}

\begin{rem}
  The contrapositive of Theorem~\ref{thm:upper} reads: a map that
  is \emph{not} a test equivalence is \emph{not} a type
  equivalence.  All type-noncontractibility results below can be
  read off from element-category computations this way; we also
  give an independent syntactic proof
  (\S\ref{sec:census}) in the tradition of \cite{CoqNote, PaperK},
  which is the only technique that works \emph{inside} the type
  structure.
\end{rem}

\section{Quotients that contract: the fixed-stratum retraction}
\label{sec:retract}

\begin{thm}
  \label{thm:retract}
  Let $h \in G_n$ generate a cyclic group $C = \langle h\rangle$
  and suppose $h$ fixes at least one cell of $\square^n$.  Then
  $\square^n/C \to 1$ is in $\Wtype$ (hence in $\Wtest$).
\end{thm}

\begin{proof}
  We give the argument for the order-$3$ element
  $h(a,b,c) = (\neg c, \neg a, b)$ of $G_3$, whose fixed cells are
  the twisted diagonal $(a, \neg a, \neg a)$; the general case is
  identical, solving around each cycle of $h$ (fixed cells exist
  precisely when every cycle has even negation weight, and the
  cycle-symmetrized meets provide invariant terms).

  Write $\sigma$ for the substitution
  $(x,y,z) \mapsto (\neg z, \neg x, y)$, so that precomposition of
  the generic cell with $\sigma$ equals the $h$-translate.  A
  homotopy $G = (G_1, G_2, G_3) : \square^3 \times \I \to
  \square^3$ descends to $\square^3/C \times \I \to \square^3/C$
  iff $G \circ \sigma = h \circ G$, that is,
  \[
    G_1 \circ \sigma = \neg G_3, \qquad
    G_2 \circ \sigma = \neg G_1, \qquad
    G_3 \circ \sigma = G_2 .
  \]
  Set $G_3 := \neg(G_1 \circ \sigma)$ and
  $G_2 := \neg(G_1 \circ \sigma^2)$: the first two equations hold
  by construction and the third follows from $\sigma^3 =
  \mathrm{id}$.  \emph{The component $G_1$ is therefore free},
  subject only to its face conditions --- and these propagate
  automatically: if $G_1|_{t=0} = x$ then $G_2|_{t=0} =
  \neg(x \circ \sigma^2) = \neg\neg y = y$ and $G_3|_{t=0} = z$;
  if $G_1|_{t=1} = m$ with $m \circ \sigma = m$, then
  $G|_{t=1} = (m, \neg m, \neg m)$, a fixed cell.  Take the
  invariant term $m = x \wedge \neg y \wedge \neg z$ and the
  standard interpolation
  \[
    G_1 = (x \wedge \neg t) \vee (m \wedge t) \vee (x \wedge m).
  \]
  All identities --- $\sigma^3 = \mathrm{id}$, invariance of $m$,
  the three equivariance equations, the six face conditions, and
  monotonicity --- are verified on all $256$ literal points
  \cite{Artifact}.  The homotopy contracts $\square^3/C$ onto its
  fixed stratum, a copy of $\I$ (the twisted diagonal embeds:
  its cells are fixed, so the quotient identifies nothing there),
  which contracts to a vertex by $a \mapsto a \vee s$ within the
  stratum.
\end{proof}

\begin{cor}
  \label{cor:squares}
  $\square^2/\langle\sw\rangle$ and
  $\square^2/\langle\gtw\rangle$ are contractible in both
  structures: for $\sw$ take the invariant term $m = x \wedge y$
  (recovering the connection contraction
  $q(s,t) \mapsto q(s \vee z, t \vee z)$), for $\gtw$ take
  $m = x \wedge \neg y$.
\end{cor}

\begin{rem}
  This is exactly what fails for cartesian cubes: Coquand's square
  is $\square^2/\langle\sw\rangle$ over a site with no
  connections, so no interpolation $G_1$ exists and the quotient
  is a genuine type-theoretic obstruction \cite{CoqNote}.  Over De
  Morgan cubes the connections dissolve every \emph{cyclic}
  fixed-stratum obstruction; what survives is the mixed case of
  \S\ref{sec:join}, where no cell is fixed by the whole group and
  no retraction target exists.
\end{rem}

\section{Free quotients}
\label{sec:free}

Call $H \leq G_n$ \emph{cell-free} if no nontrivial element fixes
a cell of $\square^n$ at any level; by the freeness theorem of
\cite{PaperK, PaperL} this holds for every subgroup generated by
elements with a cycle of odd negation weight (for $n = 2$:
$\langle\nb\rangle$, $\langle\nb_1\rangle$,
$\langle\nb_1\sw\rangle \cong \mathbb{Z}/4$,
$\langle\nb_1, \nb_2\rangle$, and conjugates).

\begin{thm}
  \label{thm:free}
  If $H$ is cell-free then $\square^n/H$ is a $K(H,1)$ in the
  cubical-type model structure.  Combined with
  $\el(\square^n/H) \simeq BH$ \cite[Lem.~elBH]{PaperL}, the two
  structures agree on the nose on all cell-free quotients.
\end{thm}

\begin{proof}[Proof sketch]
  The syntactic fibrant replacement $T(-)$ of \cite[\S4]{PaperK}
  commutes with cell-free quotients: $T(Y)/H \cong T(Y/H)$,
  because a nontrivial $h$ fixing a $\mathrm{comp}$-cell would fix
  its base data, and the base cells are free --- the induction of
  \cite[Prop.~4.3]{PaperK} generalizes verbatim from
  $\mathbb{Z}/2$ to arbitrary $H$.  Hence the fibrant replacement
  of $\square^n/H$ is $R/H$ with $R = T(\square^n)$ contractible
  and fibrant.  The projection $R \to R/H$ has unique lifting of
  open boxes (boxes are connected and the action is free), so
  $R/H$ is fibrant and the projection is a principal $H$-covering
  fibration with contractible total space: transport on the fiber
  is free and transitive, giving $\pi_1 = H$, and the higher
  homotopy of $R/H$ vanishes with that of $R$.
\end{proof}

\section{The mixed quotients: the join}
\label{sec:join}

Two subgroups of $G_2$ remain: the Klein group $\Kl$ and $G_2$
itself.  No cell of $\square^2$ is fixed by all of $\Kl$: a fixed
cell would satisfy $b = a$ (from $\sw$) and $a = \neg a$ (from
$\nb$), which is impossible.  The stabilizer of a cell $(a,b)$ is
$\langle\sw\rangle$ if $b = a$ (the \emph{diagonal stratum}),
$\langle\gtw\rangle$ if $b = \neg a$ (the \emph{antidiagonal
  stratum}), and trivial otherwise (the \emph{free stratum}).

\begin{thm}
  \label{thm:join}
  \(
    \el(\square^2/\Kl) \simeq B\mathbb{Z}/2 \star B\mathbb{Z}/2 ,
  \)
  the join of two infinite projective spaces: simply connected,
  with $\widetilde H_2 = 0$ and $H_3 = \mathbb{Z}/2$.  Moreover
  $\el(\square^2/G_2) \simeq
  \hopush(B\mathbb{Z}/2 \leftarrow B\Kl \to BG_2)$, which has
  fundamental group $\mathbb{Z}/2$.
\end{thm}

\begin{proof}
  Write $C = \el(\square^2)$, the slice $\cube/[2]$, a contractible
  category (terminal object $T = \mathrm{id}$) on which $\Kl$ acts
  by postcomposition; the category of elements of the quotient
  presheaf is the strict quotient: $\el(\square^2/\Kl) \cong
  C/\Kl$.

  \emph{Step 1 (sieve structure).}  Along any morphism
  $[f] \to [g]$ of $C/\Kl$ one has $\Stab(g) \subseteq \Stab(f)$:
  if $g \circ u = f$ and $k \in \Stab(g)$ then
  $k f = k g u = g u = f$ (the group is abelian, so stabilizers
  are well-defined on orbits).  There are no morphisms from the
  free stratum into a stabilized one, and none between the two
  stabilized strata: a morphism into the diagonal stratum forces
  equal components on a $\Kl$-orbit that has none (verified also
  by machine over all $23{,}040$ small-level composites, zero
  violations \cite{Artifact}).  So the two stabilized strata form
  a sieve $S = D \sqcup A$ and the free stratum a cosieve $U$, and
  $C/\Kl$ is the collage of the hom profunctor
  $W : S^{\mathrm{op}} \times U \to \mathrm{Set}$.

  \emph{Step 2 (the three pieces).}  The nerve of a collage is the
  double mapping cylinder
  $\hopush(NS \leftarrow N\widehat W \to NU)$, where $\widehat W$
  is the two-sided category of elements of $W$ (objects: an arrow
  from a stabilized object to a free one).
  \begin{itemize}
  \item $U$: the free stratum contains $T$ and is closed under
    maps out, so it is contractible before quotienting; the
    $\Kl$-action on it is free, so $U \simeq B\Kl$.
  \item $D$: the diagonal stratum of $C$ is
    $\el(\I) \simeq \ast$ via $a \mapsto (a,a)$; $\sw$ acts
    trivially on it and $\nb$ freely, so
    $D \simeq B(\Kl/\langle\sw\rangle) = B\mathbb{Z}/2$ ---
    a copy of $\el(L)$ \cite{PaperL}.  Likewise
    $A \simeq B(\Kl/\langle\gtw\rangle)$.
  \item $\widehat W_D$: it is isomorphic to
    $\mathrm{Arr}(D_C, U_C)/\Kl$, the quotient of the category of
    arrows from diagonal to free objects of $C$.  The projection
    $\mathrm{Arr} \to D_C$ is a Grothendieck fibration whose fiber
    over $d$ is the comma $(d \downarrow U_C)$, which has terminal
    object the unique map $d \to T$; so $\mathrm{Arr}$ is
    contractible by Quillen's Theorem~A, the $\Kl$-action on it is
    free (its free-stratum component is), and
    $\widehat W_D \simeq B\Kl$.
  \end{itemize}

  \emph{Step 3 (the maps, and the join).}  The functor
  $\widehat W_D \to U$ is the free quotient of an equivariant map
  of contractible categories, hence an equivalence $B\Kl \to
  B\Kl$; the functor $\widehat W_D \to D$ is equivariant along
  $\Kl \twoheadrightarrow \Kl/\langle\sw\rangle$, hence the
  classifying map of that projection.  Gluing the two cylinders
  along the equivalences leaves
  \[
    \el(\square^2/\Kl) \;\simeq\;
    \hopush\bigl( B(\Kl/\langle\sw\rangle) \longleftarrow B\Kl
    \longrightarrow B(\Kl/\langle\gtw\rangle) \bigr).
  \]
  Since $\Kl = \langle\sw\rangle \times \langle\gtw\rangle$, the
  two legs are the projections
  $B\mathbb{Z}/2 \times B\mathbb{Z}/2 \to B\mathbb{Z}/2$, and the
  homotopy pushout of the two projections of a product is by
  definition the join.  Joins of path-connected spaces are simply
  connected, and
  $\widetilde H_n(X \star Y) \cong \widetilde H_{n-1}(X \wedge
  Y)$ gives $\widetilde H_2 = 0$ and
  $H_3 = H_1 \otimes H_1 = \mathbb{Z}/2$.

  For $G_2$ the reflection $\nb_1$ interchanges the two stabilized
  strata, which merge into one with
  $S \simeq B(\langle \sw,\nb\rangle/\langle\sw\rangle)$, and the
  interface becomes $B\Kl$ mapping to $U \simeq BG_2$ along the
  inclusion $\Kl \hookrightarrow G_2$; van Kampen gives
  $\pi_1 = G_2/\langle\langle \sw \rangle\rangle =
  G_2/\Kl = \mathbb{Z}/2$ (the normal closure of $\sw$ contains
  its $\nb_1$-conjugate $\gtw$).
\end{proof}

\begin{cor}
  \label{cor:Wtype}
  $\square^2/\Kl \to 1 \notin \Wtype$ and
  $\square^2/G_2 \to 1 \notin \Wtype$: both quotients are
  contractible in \emph{neither} structure.
\end{cor}

\begin{proof}
  Immediate from Theorem~\ref{thm:join} and the contrapositive of
  Theorem~\ref{thm:upper}.
\end{proof}

\section{The syntactic census}
\label{sec:census}

Corollary~\ref{cor:Wtype} passes through the test structure.  We
record an independent proof of
$\square^2/\Kl \to 1 \notin \Wtype$ that never leaves the type
structure, by the factorization-invariant method of
\cite{CoqNote, PaperK}; it is the only known technique of this
kind, and it doubles as a verification of the census machinery.

Call a square $c$ of a De Morgan cubical set
\emph{$\Kl$-symmetric} if $c \circ \sigma_{\sw} = c$ and
$c \circ \sigma_{\nb} = c$, where $\sigma_{\sw}, \sigma_{\nb}$ are
the corresponding automorphisms of $[2]$ acting by precomposition
(for cells of a quotient this is orbit equality, so the condition
unfolds to a twisted equation with a twist in $\Kl$); similarly
for cubes, with the extra direction untouched.  Let $P''(X)$ be
the property: every $\Kl$-symmetric square of $X$ is connected to
a constant square by a finite chain of $\Kl$-symmetric cubes.

\begin{prop}
  \label{prop:census}
  $P''(\square^2/\Kl)$ fails: the $\Kl$-symmetric squares number
  $72$ and the chain graph decomposes into the component of the
  generic square (size $3$), the component of the constants (size
  $66$), and one further component (size $3$); the generic square
  reaches no constant.
\end{prop}

The census enumerates, for each assignment of a twist in $\Kl$ to
each of the two generators, the compatible cubes: a twist with a
transposition component determines one component of the cube from
the other over all of $\DM(3)$, and the purely diagonal twists
reduce to face-profile combinations; positive and negative
controls (the cyclic subgroups, which must connect by
Corollary~\ref{cor:squares}, and the cell-free subgroups, which
must not, by Theorem~\ref{thm:free}) all behave as theory
predicts \cite{Artifact}.  That $66$ of $72$ symmetric squares
\emph{do} contract shows the invariant is nondegenerate: the
isolation of the generic square is not a formal artifact.

\begin{prop}
  \label{prop:syntactic}
  $P''$ propagates through the syntactic fibrant replacement:
  for the factorization $1 \to T \twoheadrightarrow
  \square^2/\Kl$ of \cite[\S4]{PaperK}, $P''(1)$ implies
  $P''(T)$; and if $u : X \to Y$ has a section, $P''(X)$ implies
  $P''(Y)$.  Consequently $\square^2/\Kl \to 1 \notin \Wtype$.
\end{prop}

\begin{proof}[Proof sketch]
  As in \cite[Prop.~4.3]{PaperK}, with one new ingredient: the
  substitutions $\sigma_{\sw}, \sigma_{\nb}$ are
  \emph{automorphisms} of the free De Morgan algebra on two
  generators, so no restriction clause of a $\mathrm{comp}$-cell
  can fire under them, symmetry of a composition cell is an
  equality of $\mathrm{comp}$-constructors, and all its data
  (directions, constraint, tube, base) is itself symmetric.  The
  fresh-direction cube
  $c(x,y,t) = (v(x,y,t), \mathrm{comp}^{k \to t}(a, \psi, u,
  \langle z\rangle v(x,y,z)))$ is then $\Kl$-symmetric, agrees
  with the given square at $t = l$, and reduces the inductive
  height at $t = k$; the induction terminates at cells of the
  point.  Given $\square^2/\Kl \to 1 \in \Wtype$, two-out-of-three
  makes the vertex $1 \to \square^2/\Kl$ a weak equivalence, the
  fibration $T \to \square^2/\Kl$ is then trivial and has a
  section, and $P''(1) \Rightarrow P''(T) \Rightarrow
  P''(\square^2/\Kl)$ contradicts
  Proposition~\ref{prop:census}.
\end{proof}

\section{Agreement, non-separation, and the conjecture}
\label{sec:agree}

\begin{table}
  \centering
  \begin{tabular}{lll}
    \hline
    quotient & type structure & test structure \\
    \hline
    $L = \I/\langle\neg\rangle$
      & $K(\mathbb{Z}/2,1)$ \cite{PaperK}
      & $B\mathbb{Z}/2$ \cite{PaperL} \\
    $\square^2/\langle\sw\rangle$, $\square^2/\langle\gtw\rangle$
      & $\ast$ (Cor.~\ref{cor:squares})
      & $\ast$ (Thm.~\ref{thm:upper}) \\
    $\square^3/\mathbb{Z}/3$ twisted
      & $\ast$ (Thm.~\ref{thm:retract})
      & $\ast$ (Thm.~\ref{thm:upper}) \\
    cell-free $H$
      & $K(H,1)$ (Thm.~\ref{thm:free})
      & $BH$ \cite{PaperL} \\
    $\square^2/\Kl$
      & $\neq \ast$ (Cor.~\ref{cor:Wtype},
        Prop.~\ref{prop:syntactic})
      & $\mathbb{RP}^\infty \star \mathbb{RP}^\infty$
        (Thm.~\ref{thm:join}) \\
    $\square^2/G_2$
      & $\neq \ast$ (Cor.~\ref{cor:Wtype})
      & $\pi_1 = \mathbb{Z}/2$ (Thm.~\ref{thm:join}) \\
    \hline
  \end{tabular}
  \medskip
  \caption{Both structures agree on every tested quotient.}
  \label{tab:agree}
\end{table}

\begin{cor}
  \label{cor:nosep}
  No quotient $\square^2/H$, $H \leq G_2$, is a separating witness
  for $\Wtype \subseteq \Wtest$ on De Morgan cubical sets: the
  subgroups of $G_2$ are exhausted by
  Corollary~\ref{cor:squares} (cyclic with fixed cells),
  Theorem~\ref{thm:free} (cell-free), and
  Corollary~\ref{cor:Wtype} (mixed), and in every case the
  quotient is test-trivial if and only if it is type-trivial.
\end{cor}

\begin{conj}
  \label{conj:agree}
  The homotopy types agree on all isotropy quotients in all
  dimensions: for every $H \leq G_n$, the type structure assigns
  to $\square^n/H$ the homotopy colimit of its stabilizer strata
  --- the same diagram that computes $\el(\square^n/H)$.  In
  particular $\Wtype = \Wtest$ on De Morgan cubical sets.
\end{conj}

\begin{rem}
  \label{rem:obstacle}
  Two remarks on why the remaining inclusion resists.  First, the
  geometric realization of every quotient above is contractible
  (a folded square is a triangle), so no cellular or realization
  argument can see the answer; the phenomenon lives strictly in
  the difference between strict and homotopy quotients, on both
  sides.  Second, the element-category computation of
  Theorem~\ref{thm:join} rests on a stratification that has
  \emph{no presheaf-level counterpart}: the subpresheaf of
  $\square^2$ generated by the free cells is all of $\square^2$,
  because every diagonal cell $(a,a)$ is a face of the free cell
  $(a, a \vee (t \wedge \neg t))$ --- a Kleene-midpoint
  phenomenon with no Boolean or cartesian analogue.  A type-side
  computation of the mixed quotients therefore needs a genuinely
  new device (a presheaf-level collage), which we leave open.
\end{rem}

\begin{prop}
  \label{prop:ez}
  $\cube$ and $\cubeK$ are not Eilenberg--Zilber categories: the
  Kleene midpoint $x \wedge \neg x : [1] \to [1]$ factors through
  no (split epimorphism, monomorphism) pair.
\end{prop}

\begin{proof}
  The two points of $\I$ have the same image under
  $x \wedge \neg x$ (both go to $0 \wedge 1 = 0$), so the map is
  not a monomorphism; it is not constant either, so a
  factorization would pass through $[1]$ with a split epimorphism
  $[1] \to [1]$ whose image is the midpoint --- but a split
  epimorphism is surjective on the interval algebra, and
  $x \wedge \neg x$ generates a proper subalgebra.
\end{proof}

This explains from first principles why the elegance-based
techniques of \cite{CS} require the \emph{relative} theory over
these sites, and complements the no-go of \cite{PaperL}: on De
Morgan cubes even the absolute Eilenberg--Zilber property is
unavailable.

\subsection*{Outlook}
Three directions.  (1)~\emph{The type side of the join.}  The
natural route to Conjecture~\ref{conj:agree} for the mixed
quotients is a presheaf-level double mapping cylinder
interpolating between $\square^2/\Kl$ and the join, with the two
$L$-copies (Theorem~\ref{thm:free} and \cite{PaperK}) as the
stabilized ends.  (2)~\emph{Generation of $\Wtest$.}  The equality
$\Wtype = \Wtest$ would follow from exhibiting a generating set of
the localizer $\Wtest$ consisting of maps already in $\Wtype$;
over a totally aspherical site the candidates are exactly the
quotient maps studied here, which is what
Conjecture~\ref{conj:agree} would supply.  (3)~\emph{Beyond
  quotients.}  Corollary~\ref{cor:nosep} closes the
two-dimensional quotient route to a counterexample; a separating
witness, if any, must be either higher-dimensional with
non-abelian mixed isotropy, or not an isotropy quotient at all.

\subsection*{Acknowledgements}
The machine verifications were carried out with the toolchain of
the accompanying artifact \cite{Artifact}.

\begin{thebibliography}{99}

\bibitem{ACCRS}
S.~Awodey, E.~Cavallo, T.~Coquand, E.~Riehl, C.~Sattler,
\emph{The equivariant model structure on cartesian cubical sets},
arXiv:2406.18497.

\bibitem{BuchholtzMorehouse}
U.~Buchholtz, E.~Morehouse,
\emph{Varieties of cubical sets},
RAMiCS 2017, LNCS 10226, 77--92.

\bibitem{CS}
E.~Cavallo, C.~Sattler,
\emph{Relative elegance and cartesian cubes with one connection},
arXiv:2211.14801.

\bibitem{Cisinski}
D.-C.~Cisinski,
\emph{Les pr\'efaisceaux comme mod\`eles des types d'homotopie},
Ast\'erisque 308 (2006).

\bibitem{CCHM}
C.~Cohen, T.~Coquand, S.~Huber, A.~M\"ortberg,
\emph{Cubical type theory: a constructive interpretation of the
  univalence axiom},
TYPES 2015, LIPIcs 69, 5:1--5:34.

\bibitem{CoqNote}
T.~Coquand,
\emph{A remark on contractible families of types},
note, 2018.

\bibitem{Maltsiniotis}
G.~Maltsiniotis,
\emph{La cat\'egorie cubique avec connexions est une cat\'egorie
  test stricte},
Homology Homotopy Appl.\ 11 (2009), 309--326.

\bibitem{Sattler}
C.~Sattler,
\emph{The equivalence extension property and model structures},
arXiv:1704.06911.

\bibitem{PaperG}
R.~Shibusawa,
\emph{Boundary fibers in free De Morgan algebras},
artifact paper 6, this repository, 2026.

\bibitem{PaperK}
R.~Shibusawa,
\emph{Cubical models with reversal do not present spaces},
artifact paper 10, this repository, 2026.

\bibitem{PaperL}
R.~Shibusawa,
\emph{Reversal is unrepairable: equivariance, elegance, and the
  test comparison for De Morgan cubical sets},
artifact paper 11, this repository, 2026.

\bibitem{Artifact}
R.~Shibusawa,
\emph{Cubical strict-J artifact},
\url{https://github.com/r-shibusawa/cubical-strict-j},
release v1.12.0.

\end{thebibliography}

\end{document}
